\chapter*{Introducción}
\addcontentsline{toc}{chapter}{Introducción}

\section*{1.\quad El enfoque variacional de la mecánica.}

Desde que Newton sentó las bases sólidas de la dinámica al formular las leyes del movimiento, la ciencia de la mecánica se desarrolló a lo largo de dos líneas principales. Una rama, que llamaremos ``mecánica vectorial,''\footnote{Este uso del término no implica necesariamente que se empleen métodos vectoriales.} parte directamente de las leyes de movimiento de Newton. Su objetivo es reconocer todas las fuerzas que actúan sobre cualquier partícula dada, siendo su movimiento determinado de manera única por las fuerzas conocidas que actúan sobre ella en cada instante. El análisis y la síntesis de fuerzas y momentos es, por tanto, la preocupación básica de la mecánica vectorial.

Mientras que en la mecánica de Newton la acción de una fuerza se mide por el momento producido por dicha fuerza, el gran filósofo y universalista Leibniz, contemporáneo de Newton, propuso otra cantidad, la \textit{vis viva} (fuerza viva), como la medida adecuada para la acción dinámica de una fuerza. Esta \textit{vis viva} de Leibniz coincide ---salvo el factor inesencial 2--- con la cantidad que hoy llamamos ``energía cinética.'' Así, Leibniz reemplazó el ``momento'' de Newton por la ``energía cinética.'' Al mismo tiempo, reemplazó la ``fuerza'' de Newton por el ``trabajo de la fuerza.'' Este ``trabajo de la fuerza'' fue posteriormente reemplazado por una cantidad aún más básica, la ``función de trabajo.'' Leibniz es, por tanto, el originador de esa segunda rama de la mecánica, usualmente llamada ``mecánica analítica,''\footnote{Véase la nota sobre terminología al final del cap.\ I, sección 1.} la cual basa todo el estudio del equilibrio y el movimiento en dos cantidades escalares fundamentales, la ``energía cinética'' y la ``función de trabajo,'' esta última frecuentemente reemplazada por la ``energía potencial.''

Dado que el movimiento es, por su propia naturaleza, un fenómeno \textit{dirigido}, parece desconcertante que dos cantidades escalares deban ser suficientes para determinar el movimiento. El teorema de la energía, que establece que la suma de las energías cinética y potencial permanece invariable durante el movimiento, proporciona solamente \textit{una} ecuación, mientras que el movimiento de una sola partícula en el espacio requiere \textit{tres} ecuaciones; en el caso de sistemas mecánicos compuestos por dos o más partículas, la discrepancia se vuelve aún mayor. Y sin embargo, es un hecho que estos dos escalares fundamentales contienen la dinámica completa de incluso el sistema material más complicado, siempre que se utilicen como base de un \textit{principio} en lugar de una \textit{ecuación}.

\section*{2.\quad El procedimiento de Euler y Lagrange.}

Para ver cómo ocurre esto, consideremos una partícula que se encuentra en un punto $P_1$ en un tiempo $t_1$. Supongamos que conocemos su velocidad en ese tiempo. Supongamos además que sabemos que la partícula estará en un punto $P_2$ después de que haya transcurrido un tiempo dado. Aunque no conocemos la trayectoria seguida por la partícula, es posible establecer dicha trayectoria completamente mediante experimentación matemática, siempre que las energías cinética y potencial de la partícula estén dadas para cualquier velocidad posible y cualquier posición posible.

Euler y Lagrange, los primeros descubridores del principio exacto de mínima acción, proceden de la siguiente manera. Conectemos los dos puntos $P_1$ y $P_2$ mediante \textit{cualquier} trayectoria tentativa. Con toda probabilidad, esta trayectoria, que puede ser elegida como una curva continua arbitraria, \textit{no} coincidirá con la trayectoria real que la naturaleza ha elegido para el movimiento. Sin embargo, podemos corregir gradualmente nuestra solución tentativa y eventualmente llegar a una curva que puede ser designada como la trayectoria \textit{real} del movimiento.

Para este propósito, dejamos que la partícula se mueva a lo largo de la trayectoria tentativa de acuerdo con el principio de la energía. La suma de las energías cinética y potencial se mantiene constante y siempre igual al valor $E$ que el movimiento real ha revelado en el tiempo $t_1$. Esta restricción asigna una velocidad definida a cualquier punto de nuestra trayectoria y, por tanto, determina el movimiento. Podemos elegir libremente nuestra trayectoria, pero una vez hecho esto, la conservación de la energía determina el movimiento de manera única.

En particular, podemos calcular el tiempo en el que la partícula pasará por un punto arbitrariamente dado de nuestra trayectoria ficticia y, por tanto, la integral temporal de la \textit{vis viva}, es decir, del doble de la energía cinética, extendida sobre todo el movimiento desde $P_1$ hasta $P_2$. Esta integral temporal se denomina ``acción.'' Tiene un valor definido para nuestra trayectoria tentativa y, de igual manera, para cualquier otra trayectoria tentativa, siendo estas trayectorias siempre trazadas entre los mismos dos puntos extremos $P_1$, $P_2$ y recorridas siempre con la misma constante de energía dada $E$.

El valor de esta ``acción'' variará de una trayectoria a otra. Para algunas trayectorias resultará mayor, para otras menor. Matemáticamente, podemos imaginar que \textit{todas} las trayectorias posibles han sido probadas. Debe existir una trayectoria definida (al menos si $P_1$ y $P_2$ no están demasiado alejados) para la cual la acción asume un valor mínimo. El principio de mínima acción afirma que \textit{esta trayectoria particular es la elegida por la naturaleza como la trayectoria real del movimiento}.

Hemos explicado la operación del principio para una sola partícula. Puede generalizarse, sin embargo, a cualquier número de partículas y a cualquier sistema mecánico arbitrariamente complicado.

\section*{3.\quad El procedimiento de Hamilton.}

Encontramos problemas de mecánica para los cuales la función de trabajo es una función no solo de la posición de la partícula sino también del tiempo. Para tales sistemas, la ley de conservación de la energía no se cumple, y el principio de Euler y Lagrange no es aplicable, pero sí lo es el de Hamilton.

En el procedimiento de Hamilton, comenzamos nuevamente con el punto inicial dado $P_1$ y el punto final dado $P_2$. Pero ahora no restringimos el movimiento de prueba de ninguna manera. No solo la trayectoria puede ser elegida arbitrariamente ---salvo las condiciones naturales de continuidad--- sino que también el movimiento en el tiempo está a nuestra disposición. Todo lo que requerimos ahora es que nuestro movimiento tentativo comience en el tiempo observado $t_1$ del movimiento real y termine en el tiempo observado $t_2$. (Esta condición no se satisface en el procedimiento de Euler--Lagrange, porque allí el teorema de la energía restringe el movimiento, y el tiempo empleado para ir de $P_1$ a $P_2$ en el movimiento tentativo generalmente diferirá del tiempo empleado en el movimiento real.)

La cantidad característica que ahora usamos como medida de la acción ---no existe desafortunadamente un nombre estándar adoptado para esta cantidad--- es la integral temporal de la \textit{diferencia entre las energías cinética y potencial}. La formulación hamiltoniana del principio de mínima acción afirma que \textit{el movimiento real realizado en la naturaleza es aquel movimiento particular para el cual esta acción asume su valor mínimo}.

Se puede demostrar que, en el caso de sistemas ``conservativos,'' es decir, sistemas que satisfacen la ley de conservación de la energía, el principio de Euler--Lagrange es una consecuencia del principio de Hamilton, pero este último principio permanece válido incluso para sistemas no conservativos.

\section*{4.\quad El cálculo de variaciones.}

El problema matemático de minimizar una integral se trata dentro de una rama especial del cálculo, llamada ``cálculo de variaciones.'' La teoría matemática muestra que nuestros resultados finales pueden ser establecidos sin tener en cuenta la infinidad de trayectorias tentativamente posibles. Podemos restringir nuestro experimento matemático a aquellas trayectorias que estén \textit{infinitamente cerca} de la trayectoria real. Una trayectoria tentativa que difiere de la trayectoria real en un grado arbitrario pero aun así \textit{infinitesimal}, se denomina una ``variación'' de la trayectoria real, y el cálculo de variaciones investiga los cambios en el valor de una integral causados por tales variaciones infinitesimales de la trayectoria.

\section*{5.\quad Comparación entre los tratamientos vectorial y variacional de la mecánica.}

Los tratamientos vectorial y variacional de la mecánica son dos descripciones matemáticas diferentes del mismo ámbito de fenómenos naturales. La teoría de Newton se basa en dos vectores fundamentales: ``momento'' y ``fuerza''; la teoría variacional, fundada por Euler y Lagrange, basa todo en dos cantidades escalares: ``energía cinética'' y ``función de trabajo.'' Aparte de la conveniencia matemática, puede plantearse la cuestión de la equivalencia de estas dos teorías. En el caso de partículas \textit{libres}, es decir, partículas cuyo movimiento no está restringido por ``ligaduras'' dadas, las dos formas de descripción conducen a resultados equivalentes. Pero para sistemas con ligaduras, el tratamiento analítico es más simple y más económico. Las ligaduras dadas se toman en cuenta de manera natural al dejar que el sistema se mueva a lo largo de todas las trayectorias tentativas en armonía con ellas. El tratamiento vectorial tiene que dar cuenta de las fuerzas que mantienen las ligaduras y tiene que hacer hipótesis definidas acerca de ellas. La tercera ley de movimiento de Newton, ``acción igual a reacción,'' no abarca todos los casos. Es suficiente solo para la dinámica de cuerpos rígidos.

Por otro lado, el enfoque de Newton no restringe la naturaleza de una fuerza, mientras que el enfoque variacional supone que las fuerzas que actúan son derivables de una cantidad escalar, la ``función de trabajo.'' Las fuerzas de naturaleza friccional, que no tienen función de trabajo, están fuera del ámbito de los principios variacionales, mientras que el esquema newtoniano no tiene dificultad en incluirlas.

\begin{small}
\noindent Tales fuerzas se originan a partir de fenómenos intermoleculares que son despreciados en la descripción macroscópica del movimiento. Si los parámetros macroscópicos de un sistema mecánico se completan con la adición de parámetros microscópicos, fuerzas no derivables de una función de trabajo con toda probabilidad no aparecerían.
\end{small}

\section*{6.\quad Evaluación matemática de los principios variacionales.}

Muchos problemas elementales de física e ingeniería son resolubles mediante mecánica vectorial y no requieren la aplicación de métodos variacionales. Pero en todos los problemas más complicados, la superioridad del tratamiento variacional se hace conspicua. Esta superioridad se debe a la completa libertad que tenemos para elegir las coordenadas apropiadas para nuestro problema. Los problemas que son bien adecuados para el tratamiento vectorial son esencialmente aquellos que pueden manejarse con un sistema de referencia rectangular, ya que la descomposición de vectores en coordenadas curvilíneas es un procedimiento engorroso si no se guía por los principios avanzados del cálculo tensorial. Aunque la importancia fundamental de los invariantes y covariantes para todos los fenómenos de la naturaleza ha sido descubierta solo recientemente y, por tanto, no era conocida en la época de Euler y Lagrange, el enfoque variacional de la mecánica logró anticipar este desarrollo al satisfacer el principio de invariancia automáticamente. Se nos permite soberana libertad para elegir nuestras coordenadas, ya que nuestros procesos y las ecuaciones resultantes permanecen válidos para una elección arbitraria de coordenadas. El valor matemático y filosófico del método variacional está firmemente anclado en esta libertad de elección y en la correspondiente libertad de transformaciones de coordenadas arbitrarias. Facilita enormemente la formulación de las ecuaciones diferenciales del movimiento, así como su solución. Si damos con cierto tipo de coordenadas, llamadas ``cíclicas'' o ``ignorables,'' se logra de inmediato una integración parcial de las ecuaciones diferenciales básicas. Si \textit{todas} nuestras coordenadas son ignorables, nuestro problema está completamente resuelto. De ahí que podamos formular el problema completo de resolver las ecuaciones diferenciales del movimiento como un problema de transformación de coordenadas. En lugar de intentar integrar directamente las ecuaciones diferenciales del movimiento, intentamos producir más y más coordenadas ignorables. En la forma de Euler--Lagrange de la mecánica, es más o menos accidental si damos con las coordenadas correctas, porque no tenemos una manera sistemática de producir coordenadas ignorables. Pero los desarrollos posteriores de la teoría por Hamilton y Jacobi ampliaron enormemente los procedimientos originales al introducir las ``ecuaciones canónicas,'' con sus propiedades de transformación mucho más amplias. Aquí somos capaces de producir un conjunto \textit{completo} de coordenadas ignorables resolviendo una sola ecuación diferencial parcial.

Aunque la solución real de esta ecuación diferencial es posible solo para una clase restringida de problemas, sucede que muchos problemas importantes de la física teórica pertenecen precisamente a esta clase. Y así, la forma más avanzada de la mecánica analítica resulta ser no solo estéticamente y lógicamente la más satisfactoria, sino al mismo tiempo muy práctica, al proporcionar una herramienta para la solución de muchos problemas dinámicos que no son accesibles a métodos elementales.\footnote{El presente libro no discute otros métodos de integración que no están basados en la teoría de transformaciones. Respecto a tales métodos, se remite al lector a los libros de texto avanzados mencionados en la Bibliografía.}

\section*{7.\quad Evaluación filosófica del enfoque variacional de la mecánica.}

Aunque se acepta tácitamente hoy en día que los tratados científicos deben evitar las discusiones filosóficas, en el caso de los principios variacionales de la mecánica puede tolerarse una excepción a la regla, en parte porque estos principios están enraizados en una época que estaba filosóficamente orientada en un grado muy alto, y en parte porque el método variacional ha sido frecuentemente el foco de controversias filosóficas y malinterpretaciones.

En efecto, la idea de ampliar la realidad al incluir posibilidades ``tentativas'' y luego seleccionar una de ellas mediante la condición de que minimiza cierta cantidad, parece traer un \textit{propósito} al flujo de los eventos naturales. Esto está en contradicción con la descripción \textit{causal} usual de las cosas. Sin embargo, no debemos sorprendernos de que, para el enfoque más universal que era corriente en los siglos XVII y XVIII, las dos formas de pensar no aparecieran necesariamente como contradictorias. La nota dominante de todo ese período era la armonía aparentemente preestablecida entre ``razón'' y ``mundo.'' Los grandes descubrimientos de las matemáticas superiores y su aplicación inmediata a la naturaleza imbuían a los filósofos de aquellos días con una confianza ilimitada en la estructura fundamentalmente intelectual del mundo. El ``deus intellectualis'' era el tema básico de la filosofía de Leibniz, no menos que la de Spinoza. Al mismo tiempo, Leibniz tenía fuertes tendencias teleológicas, y sus actividades tuvieron no poca influencia en el desarrollo de los métodos variacionales.\footnote{Véase el atractivo estudio histórico de A.\ Kneser, \textit{Das Prinzip der kleinsten Wirkung von Leibniz bis zur Gegenwart} (Leipzig: Teubner, 1928).} Pero esto no es sorprendente si observamos cómo el interés puramente estético y lógico en los problemas de máximos y mínimos dio uno de los impulsos más fuertes para el desarrollo del cálculo infinitesimal, y cómo la derivación por Fermat de las leyes de la óptica geométrica sobre la base de su ``principio de la llegada más rápida'' no podía dejar de impresionar a los científicos filosóficamente orientados de aquellos días. Que el uso diletante de estas tendencias por Maupertuis y otros con fines teológicos haya puesto toda la tendencia en descrédito, no es culpa de los grandes filósofos.

El sobrio y práctico siglo XIX, orientado hacia los hechos ---que se extiende hasta nuestros días--- sospechaba de todas las tendencias especulativas e interpretativas como ``metafísicas'' y limitaba su programa a la pura descripción de los eventos naturales. En esta filosofía, las matemáticas desempeñan el papel de un método abreviado, un lenguaje convenientemente económico para expresar relaciones involucradas. De ahí que no sea sorprendente, sino bastante consistente con el espíritu ``positivista'' del siglo XIX, encontrar la siguiente valoración de la mecánica analítica por parte de una de las principales figuras de esa tendencia, E.\ Mach, en ``La Ciencia de la Mecánica'' (\textit{Open Court}, 1893, p.\ 480): ``No cabe esperar ninguna luz fundamental de esta rama de la mecánica. Por el contrario, el descubrimiento de cuestiones de principio debe estar sustancialmente completado antes de que podamos pensar en enmarcar la mecánica analítica, cuyo único objetivo es un dominio \textit{práctico} perfecto de los problemas. Quienquiera que confunda esta situación nunca comprenderá la gran obra de Lagrange, cuyo desempeño aquí también es esencialmente de carácter \textit{económico}.'' (Cursivas en el original.) Según esta filosofía, los principios variacionales de la mecánica no son más que formulaciones matemáticas alternativas de las leyes fundamentales de Newton, sin ninguna importancia primaria.

Sin embargo, las tendencias filosóficas van y vienen, y la última palabra nunca está dicha. En nuestro propio tiempo hemos presenciado al menos un descubrimiento fundamental de magnitud sin precedentes, a saber, la Teoría de la Relatividad General de Einstein, que fue obtenida mediante especulación matemática y filosófica del más alto orden. Aquí se realizó un descubrimiento mediante un tipo de razonamiento que un positivista no puede dejar de llamar ``metafísico,'' y sin embargo proporcionó una visión profunda del corazón de las cosas que la mera experimentación y el registro sobrio de hechos nunca podrían haber revelado. La Teoría de la Relatividad General trajo nuevamente a primer plano el espíritu de los grandes teóricos cósmicos de Grecia y del siglo XVIII.

A la luz de los descubrimientos de la relatividad, el fundamento variacional de la mecánica merece más que una valoración puramente formalista. Lejos de ser nada más que una formulación alternativa de las leyes de movimiento de Newton, los siguientes puntos sugieren la supremacía del método variacional:

\begin{enumerate}
\item El Principio de Relatividad requiere que las leyes de la naturaleza se formulen de manera ``invariante,'' es decir, independientemente de cualquier sistema de referencia especial. Los métodos del cálculo de variaciones satisfacen automáticamente este principio, porque el mínimo de una cantidad escalar no depende de las coordenadas en las cuales dicha cantidad se mide. Mientras que las ecuaciones de movimiento de Newton no satisfacían el principio de relatividad, el principio de mínima acción permaneció válido, con la única modificación de que la cantidad de acción básica debía ponerse en armonía con el requisito de invariancia.

\item La Teoría de la Relatividad General ha demostrado que la materia no puede separarse del campo y es de hecho una extensión del campo. De ahí que las ecuaciones básicas de la física deban formularse como ecuaciones diferenciales parciales en lugar de ecuaciones diferenciales ordinarias. Mientras que la imagen de partículas de Newton difícilmente puede ponerse en armonía con el concepto de campo, los métodos variacionales no están restringidos a la mecánica de partículas sino que pueden extenderse a la mecánica de continuos.

\item El Principio de la Relatividad General se satisface automáticamente si la ``acción'' fundamental del principio variacional es elegida como un invariante bajo cualquier transformación de coordenadas. Puesto que la geometría diferencial de Riemann nos proporciona tales invariantes, no tenemos dificultad en establecer las ecuaciones de campo requeridas. Aparte de esto, nuestro conocimiento actual de las matemáticas no nos da ninguna pista para la formulación de un sistema covariante, y al mismo tiempo consistente, de ecuaciones de campo. De ahí que, a la luz de la relatividad, la aplicación del cálculo de variaciones a las leyes de la naturaleza asuma más que una significación accidental.
\end{enumerate}
